% part: intuitionistic-logic
% chapter: introduction
% section: syntax

\documentclass[../../../include/open-logic-section]{subfiles}

\begin{document}

\olfileid{int}{int}{syn}

\olsection{Sintaks Logika Intuisionistik}

Sintaks logika intuisionistik padha karo sintaks logika proposisional.
Ing logika proposisional klasik, sawetara konektif bisa ditetepake
nganggo konektif liyane. Umpamane, $!A \lif !B$ bisa ditetepake
nganggo $\lnot !A \lor !B$, utawa $!A \lor !B$ nganggo
$\lnot(\lnot !A \land \lnot
!B)$. Mula, andharan logika klasik kerep ngenalake sawetara
konektif minangka cekakan saka definisi mau. Ing logika
intuisionistik iki ora lumaku, kajaba ana rong prakara: $\lnot !A$
bisa---lan kerep---ditetepake minangka cekakan saka $!A \lif \lfalse$.
Yen mangkono, $\lfalse$ ora kena ditetepake minangka cekakan!{} Kajaba
iku, $!A \liff !B$ bisa ditetepake kaya ing logika klasik minangka
$(!A \lif !B) \land (!B \lif
!A)$.

!!^{formula}s logika proposisional intuisionistik dibangun saka
\emph{!!{propositional variable}s} lan konstanta
proposisional~$\lfalse$ nganggo \emph{konektif logis}. Unsur-unsure yaiku:

\begin{enumerate}
\item Himpunan~$\PVar$ kang !!^a{denumerable} saka !!{propositional variable}s $\Obj p_0$,
  $\Obj p_1$, \dots
  \item Konstanta proposisional kanggo !!{falsity}~$\lfalse$.
\item Konektif logis: $\land$
  (konjungsi), $\lor$ (disjungsi), $\lif$ (kondisional).
\item Tandha wacan: (, ), lan koma.
\end{enumerate}

\begin{defn}[Rumus]
\ollabel{defn:formulas} Himpunan~$\Frm[L_0]$ saka \emph{!!{formula}s}
logika proposisional intuisionistik ditetepake kanthi induksi kaya mangkene:
\begin{enumerate}
\item $\lfalse$ iku !!{formula} atomik.

\item Saben !!{propositional variable}~$\Obj p_i$ iku !!{formula}
  atomik.

\item Yen $!A$ lan $!B$ iku !!{formula}s, $(!A \land
  !B)$ uga !!a{formula}.

\item Yen $!A$ lan $!B$ iku !!{formula}s, $(!A \lor !B)$
  uga !!a{formula}.

\item Yen $!A$ lan $!B$ iku !!{formula}s, $(!A \lif !B)$
  uga !!a{formula}.

\tagitem{limitClause}{Ora ana liyane kang dadi !!a{formula}.}{}
\end{enumerate}
\end{defn}

Saliyane konektif primitif ing ndhuwur, kita uga nggunakake lambang
kang \emph{ditetepake} iki: $\lnot$ (negasi) lan $\liff$
(!!{biconditional}). Rumus kang dibangun nganggo operator kang
ditetepake iku kudu dimangerteni kaya mangkene:
\begin{enumerate}
\item $\lnot !A$ iku cekakan saka $!A \lif \lfalse$.

\item $!A \liff !B$ iku cekakan saka $(!A \lif !B) \land (!B
  \lif !A)$.
\end{enumerate}

Sanadyan $\lnot$ kanthi resmi dianggep cekakan, kadhangkala kita
menehi aturan lan klausa definisi kang cetha kanggo~$\lnot$
kaya-kaya iku konektif primitif. Tujuane utamane supaya kita bisa
menehi soal latihan.

\end{document}
